% ===================================================================
%  Семинар 4. Содержательная часть — единый источник для обоих файлов.
%
%  Состав и порядок задач — по «Семинар 4 (дистанционно).docx»:
%      4.1         — несколько структурных сдвигов (фиктивные переменные);
%      10.4        — INV, GNP, CONS: диагностика мультиколлинеарности;
%      4.17, 4.18  — функция Кобба–Дугласа: мультиколлинеарность и
%                    ограничение постоянной отдачи на масштаб;
%      4.1.1       — строгая мультиколлинеарность и VIF по матрице X.
%  Формулы VIF и F_j из того же файла вынесены в раздел «Основные формулы».
%
%  Источники:
%      4.1, 4.17, 4.18 — Сборник задач к начальному курсу эконометрики
%                        (РАНХиГС), гл. 4, с. 104, 128–130;
%      10.4            — Бородич С. А. Вводный курс эконометрики, гл. 10
%                        «Мультиколлинеарность», упражнения и задачи, № 4;
%      4.1.1           — Задачник ВШЭ, § 4.1 «Регрессионная диагностика».
%
%  Раздел «Основные формулы» и всё внутри \begin{solution}...\end{solution}
%  попадает только в seminar_04_solutions.pdf. Блоки \begin{excel}...\end{excel}
%  печатаются в обоих файлах: это часть задания.
%
%  Числовые результаты пересчитаны независимо. Задача 10.4 совпадает с
%  выводом Пакета анализа в «Семинар 4.xlsx», задачи 4.17 и 4.18 — с
%  таблицами на с. 129–130 сборника.
% ===================================================================

\seminartitle{Семинар 4}{Структурные сдвиги и мультиколлинеарность}

% ===================================================================
% Теория печатается только в файле с решениями.
\ifsolutions
\section{Основные формулы}

% -------------------------------------------------------------------
\subsection{Модель и основные тесты}\label{sec:model}

Модель множественной регрессии в матричной форме:
\begin{equation}
  \mathbf{y} = \mathbf{X}\boldsymbol{\beta} + \boldsymbol{\varepsilon},
  \qquad
  \widehat{\boldsymbol{\beta}} = (\mathbf{X}^{\top}\mathbf{X})^{-1}\mathbf{X}^{\top}\mathbf{y},
  \qquad
  \D(\widehat{\boldsymbol{\beta}}) = \sigma^2(\mathbf{X}^{\top}\mathbf{X})^{-1},
  \label{eq:ols}
\end{equation}
где $\mathbf{X}$ — матрица регрессоров размера $n\times k$, $k$ — число
параметров \emph{вместе} со свободным членом. Оценка существует только тогда,
когда столбцы $\mathbf{X}$ линейно независимы: иначе матрица
$\mathbf{X}^{\top}\mathbf{X}$ вырождена.

Отдельный коэффициент проверяется $t$-тестом, а значимость регрессии в
целом — $F$-тестом:
\begin{equation}
  t = \frac{\widehat{\beta}_j}{s_{\widehat{\beta}_j}} \sim t(n-k),
  \qquad
  F = \frac{\RSS/(k-1)}{\ESS/(n-k)}
    = \frac{R^2}{1-R^2}\cdot\frac{n-k}{k-1} \sim F(k-1,\ n-k).
  \label{eq:tf}
\end{equation}
Как и раньше, $\RSS$ — объяснённая сумма квадратов, $\ESS = \sum e_t^2$ —
остаточная, $R^2 = \RSS/\TSS = 1 - \ESS/\TSS$. Модели с разным числом
регрессоров сравнивают по скорректированному коэффициенту детерминации
\begin{equation}
  R^2_{\text{adj}}
  = 1 - \frac{\ESS/(n-k)}{\TSS/(n-1)}
  = 1 - \frac{n-1}{n-k}\bigl(1 - R^2\bigr)
  = R^2 - \frac{k-1}{n-k}\bigl(1 - R^2\bigr).
  \label{eq:r2adj}
\end{equation}
Первая запись — определение: каждая сумма квадратов делится на свои степени
свободы. Вторая получается из неё, если учесть $\ESS/\TSS = 1 - R^2$. Третья —
раскрытием скобок: она показывает, что $R^2_{\text{adj}} \le R^2$ (равенство
только при $k = 1$ или $R^2 = 1$), что поправка тем больше, чем больше
параметров $k$ при том же $n$, и что $R^2_{\text{adj}}$ может быть
отрицательным.

% -------------------------------------------------------------------
\subsection{Тест на линейные ограничения}\label{sec:restr}

Пусть на коэффициенты наложено $q$ линейных ограничений. Модель оценивают
без ограничений (UR) и с ограничениями (R); во второй ограничения подставлены
в уравнение, так что оцениваемых параметров становится меньше. Всегда
$\ESS_{\mathrm{R}} \ge \ESS_{\mathrm{UR}}$, и вопрос в том, велик ли прирост:
\begin{equation}
  F = \frac{\bigl(\ESS_{\mathrm{R}} - \ESS_{\mathrm{UR}}\bigr)/q}
           {\ESS_{\mathrm{UR}}/(n-k)} \sim F(q,\ n-k),
  \label{eq:restr}
\end{equation}
где $k$ — число параметров модели \emph{без} ограничений. Если у обеих моделей
одна и та же зависимая переменная, то и $\TSS$ одна и та же, и формулу можно
переписать через $R^2$:
\[
  F = \frac{\bigl(R^2_{\mathrm{UR}} - R^2_{\mathrm{R}}\bigr)/q}
           {\bigl(1 - R^2_{\mathrm{UR}}\bigr)/(n-k)} .
\]

\begin{remark}
Форма через $R^2$ верна \emph{только} при общей зависимой переменной. Если
после подстановки ограничения в левой части стоит другая величина (например,
$\ln(Q/K)$ вместо $\ln Q$, задача~4.18), то $\TSS$ у моделей разные, их $R^2$
несравнимы, и пользоваться надо только формой \eqref{eq:restr} через $\ESS$:
остатки у обеих моделей измеряют ошибку в одной и той же величине.
\end{remark}

% -------------------------------------------------------------------
\subsection{Фиктивные переменные и структурные сдвиги}\label{sec:dummy}

Фиктивная переменная $D_t \in \{0, 1\}$ позволяет сдвинуть свободный член и
изменить наклон:
\begin{equation}
  y_t = \alpha + \beta x_t + \gamma D_t + \delta\,(D_tx_t) + \varepsilon_t :
  \qquad
  \begin{cases}
    y = \alpha + \beta x, & D_t = 0, \\[2pt]
    y = (\alpha + \gamma) + (\beta + \delta)\,x, & D_t = 1 .
  \end{cases}
  \label{eq:dummy}
\end{equation}
Если структурный сдвиг произошёл в момент $t_0$, берут $D_t = 1$ при
$t > t_0$ и $D_t = 0$ иначе. Отсутствие сдвига — гипотеза
$H_0\colon \gamma = \delta = 0$, $q = 2$ ограничения в \eqref{eq:restr}.

\emph{Несколько сдвигов.} Для каждого момента $t_j$ вводится своя пара
«фиктивная переменная и её произведение на $x$». Если выборка разбита на $p$
частей и в каждой оцениваются все $k$ параметров, модель с фиктивными
переменными подгоняет части независимо, поэтому
$\ESS_{\mathrm{UR}} = \ESS_1 + \dots + \ESS_p$. Тест Чоу для $p$ частей:
\begin{equation}
  F = \frac{\bigl(\ESS_{\mathrm{R}} - (\ESS_1 + \dots + \ESS_p)\bigr)\big/\bigl((p-1)k\bigr)}
           {(\ESS_1 + \dots + \ESS_p)\big/(n - pk)}
  \sim F\bigl((p-1)k,\ n - pk\bigr),
  \label{eq:chow}
\end{equation}
где $\ESS_{\mathrm{R}}$ — остаточная сумма одной регрессии по всей выборке.
При $p = 2$ это обычный тест Чоу.

\emph{Излом без скачка.} Если функция должна оставаться непрерывной, а
меняется только наклон, используют переменную $(x_t - x^{*})D_t$: в точке
$x^{*}$ добавка обращается в нуль, и скачка нет.

% -------------------------------------------------------------------
\subsection{Мультиколлинеарность}\label{sec:mc}

\emph{Строгая (точная) мультиколлинеарность} — линейная зависимость столбцов
$\mathbf{X}$: найдутся числа $c_1, \dots, c_k$, не все нулевые, такие что
$c_1\mathbf{x}_1 + \dots + c_k\mathbf{x}_k = \mathbf{0}$. Тогда
$\operatorname{rank}\mathbf{X} < k$, $\det(\mathbf{X}^{\top}\mathbf{X}) = 0$, и
оценка \eqref{eq:ols} не существует: разные наборы коэффициентов дают один и
тот же прогноз, и данные не могут выбрать между ними.

\emph{Нестрогая мультиколлинеарность} — сильная, но не точная связь
регрессоров. Оценки существуют и остаются несмещёнными и эффективными
(теорема Гаусса--Маркова не нарушена), но:
\begin{itemize}[leftmargin=2.2em, itemsep=2pt]
  \item дисперсии оценок велики, и отдельные $t$-статистики малы, хотя $R^2$
        высокий и $F$-тест \eqref{eq:tf} уверенно отвергает незначимость
        регрессии в целом;
  \item оценки неустойчивы: добавление или удаление нескольких наблюдений
        сильно меняет коэффициенты;
  \item знаки коэффициентов могут противоречить содержательному смыслу.
\end{itemize}
Причина в том, что коэффициент $\beta_j$ показывает эффект изменения $x_j$
\emph{при неизменных} остальных регрессорах, а в данных такие изменения почти
не встречаются: регрессоры движутся вместе. Информации для разделения их
вкладов мало.

\textbf{Диагностика.}
\begin{enumerate}[label=\arabic*), leftmargin=2.2em, itemsep=2pt]
  \item Парные коэффициенты корреляции регрессоров: $|r| > 0{,}8$ —
        тревожный признак. Но при трёх и более регрессорах линейная связь
        может быть сильной и при умеренных парных корреляциях.
  \item \emph{Фактор инфляции дисперсии.} Для $j$-го регрессора оценивают
        вспомогательную регрессию $x_j$ на все остальные регрессоры (с
        константой) и берут её коэффициент детерминации $R_j^2$:
        \begin{equation}
          \mathrm{VIF}_j = \frac{1}{1 - R_j^2},
          \qquad
          \D(\widehat{\beta}_j)
          = \mathrm{VIF}_j\cdot\frac{\sigma^2}{\sum_t (x_{tj} - \overline{x}_j)^2} .
          \label{eq:vif}
        \end{equation}
        $\mathrm{VIF}_j$ показывает, во сколько раз дисперсия оценки больше,
        чем была бы при регрессоре, не связанном с остальными. Правило:
        $\mathrm{VIF}_j > 10$ (то есть $R_j^2 > 0{,}9$) — серьёзная
        мультиколлинеарность. VIF зависит только от матрицы $\mathbf{X}$, но
        не от $\mathbf{y}$. При двух регрессорах $R_j^2 = r^2(x_1, x_2)$ и
        оба VIF одинаковы.
  \item \emph{$F$-тест вспомогательной регрессии.} Если в модели $m$
        объясняющих переменных (без константы), то во вспомогательной
        регрессии $x_j$ на остальные $m - 1$ переменных и константу $m$
        параметров, и её значимость в целом проверяется по \eqref{eq:tf}:
        \begin{equation}
          F_j = \frac{R_j^2}{1 - R_j^2}\cdot\frac{n-m}{m-1} \sim F(m-1,\ n-m).
          \label{eq:fj}
        \end{equation}
        Если $F_j > F_{0{,}05}(m-1, n-m)$, регрессор $x_j$ значимо линейно
        связан с остальными.
\end{enumerate}

% -------------------------------------------------------------------
\subsection{Короткая и длинная регрессии}\label{sec:short}

Пусть оценены «длинная» регрессия $y$ на $x_1$ и $x_2$ и «короткая» — $y$
только на $x_1$ (обе с константой), а $\widehat{\delta}$ — наклон
вспомогательной регрессии $x_2$ на $x_1$. Тогда для МНК-оценок выполняется
алгебраическое тождество
\begin{equation}
  \widetilde{\beta}_1 = \widehat{\beta}_1 + \widehat{\beta}_2\,\widehat{\delta},
  \label{eq:short}
\end{equation}
где $\widetilde{\beta}_1$ — наклон короткой регрессии, $\widehat{\beta}_1$,
$\widehat{\beta}_2$ — оценки длинной. Смысл: в короткой регрессии рост $x_1$
на единицу сопровождается в среднем ростом $x_2$ на $\widehat{\delta}$, и
коэффициент $\widetilde{\beta}_1$ собирает \emph{оба} эффекта — прямой
$\widehat{\beta}_1$ и косвенный $\widehat{\beta}_2\widehat{\delta}$. Поэтому
при сильно связанных регрессорах короткая регрессия хорошо прогнозирует, но
её коэффициент нельзя толковать как эффект одного $x_1$.

% -------------------------------------------------------------------
\subsection{Способы борьбы с мультиколлинеарностью}\label{sec:cure}

\begin{enumerate}[label=\arabic*), leftmargin=2.2em, itemsep=2pt]
  \item \emph{Ничего не делать}, если цель — прогноз: при сохранении той же
        связи регрессоров прогноз остаётся точным.
  \item \emph{Исключить} одну из связанных переменных. Цена — смещение
        оставшихся коэффициентов \eqref{eq:short}, если исключённая переменная
        на самом деле влияет на $y$.
  \item \emph{Использовать априорную информацию}: наложить на коэффициенты
        ограничение, известное из теории (например, $\beta_1 + \beta_2 = 1$), и
        подставить его в модель. Число оцениваемых параметров уменьшается.
        Ограничение стоит предварительно проверить по \eqref{eq:restr}.
  \item \emph{Преобразовать переменные}: перейти к разностям, отношениям
        (на душу населения, к ВВП), к комбинациям с экономическим смыслом.
  \item \emph{Увеличить выборку} или найти данные, где регрессоры меняются
        независимо друг от друга.
\end{enumerate}

% -------------------------------------------------------------------
\subsection{Производственная функция Кобба--Дугласа}\label{sec:cd}

Функция $Q = \alpha L^{\beta_1}K^{\beta_2}$ нелинейна, но после
логарифмирования становится линейной по параметрам:
\begin{equation}
  \ln Q = \ln\alpha + \beta_1\ln L + \beta_2\ln K + \varepsilon .
  \label{eq:cd}
\end{equation}
Коэффициенты — эластичности: $\beta_1 = \partial\ln Q/\partial\ln L$ —
на сколько процентов растёт выпуск при росте труда на 1\,\%. Сумма
$\beta_1 + \beta_2$ — отдача от масштаба: при увеличении обоих факторов в
$\lambda$ раз выпуск растёт в $\lambda^{\beta_1+\beta_2}$ раз. Постоянная отдача
— $\beta_1 + \beta_2 = 1$.
\fi

% ===================================================================
\section{Задачи}

% -------------------------------------------------------------------
\problem{4.1}

С помощью бинарных переменных напишите уравнение, соответствующее наличию
двух структурных изменений в моменты времени $t_0$ и $t_1$ (предполагается, что
$t_0 < t_1$).

\begin{solution}
Пусть $y_t$ — зависимая переменная, $x_t$ — регрессор, наблюдения — временные
ряды, $t = 1, \dots, n$. Два структурных изменения делят выборку на три
промежутка: $t \leqslant t_0$, $t_0 < t \leqslant t_1$ и $t > t_1$.

\medskip
\textbf{1) Модель.} Вводим по фиктивной переменной на каждый сдвиг:
\[
  R_t = \begin{cases} 0, & t \leqslant t_0, \\ 1, & t > t_0, \end{cases}
  \qquad
  S_t = \begin{cases} 0, & t \leqslant t_1, \\ 1, & t > t_1, \end{cases}
\]
и включаем в регрессию сами переменные и их произведения на $x_t$
\eqref{eq:dummy}:
\[
  y_t = \beta_0 + \beta_1x_t
        + \gamma_1R_t + \delta_1(R_tx_t)
        + \gamma_2S_t + \delta_2(S_tx_t) + \varepsilon_t .
\]
Здесь $\gamma_1$, $\gamma_2$ сдвигают свободный член, а $\delta_1$, $\delta_2$
меняют наклон. Переменная $S_t$ включается \emph{поверх} $R_t$, поэтому
$\gamma_2$ и $\delta_2$ — изменения в момент $t_1$ по сравнению со вторым
промежутком, а не с первым.

\medskip
\textbf{2) Уравнения на промежутках.} Подставляем значения $R_t$ и $S_t$:
\begin{center}
\small
\setlength{\tabcolsep}{4pt}
\begin{tabular}{lcc>{$}l<{$}>{$}c<{$}}
\toprule
Промежуток & $R_t$ & $S_t$ & \text{Уравнение на промежутке} & \text{Наклон} \\
\midrule
$t \leqslant t_0$       & 0 & 0 & y = \beta_0 + \beta_1x & \beta_1 \\
$t_0 < t \leqslant t_1$ & 1 & 0 & y = (\beta_0 + \gamma_1) + (\beta_1 + \delta_1)\,x & \beta_1 + \delta_1 \\
$t > t_1$               & 1 & 1 & y = (\beta_0 + \gamma_1 + \gamma_2) + (\beta_1 + \delta_1 + \delta_2)\,x & \beta_1 + \delta_1 + \delta_2 \\
\bottomrule
\end{tabular}
\end{center}

\textbf{Скачки.} Скачок — разность значений \emph{новой} и \emph{старой}
прямой в момент сдвига. Обе прямые берём из таблицы выше и подставляем в них
одно и то же значение регрессора — то, что было в момент сдвига:
\begin{center}
\small
\setlength{\tabcolsep}{4pt}
\begin{tabular}{c>{$}l<{$}>{$}l<{$}>{$}c<{$}}
\toprule
Сдвиг & \text{Старая прямая в точке сдвига} & \text{Новая прямая в точке сдвига} & \text{Скачок} \\
\midrule
$t_0$ & \beta_0 + \beta_1x_{t_0}
      & (\beta_0 + \gamma_1) + (\beta_1 + \delta_1)\,x_{t_0}
      & \gamma_1 + \delta_1x_{t_0} \\
$t_1$ & (\beta_0 + \gamma_1) + (\beta_1 + \delta_1)\,x_{t_1}
      & (\beta_0 + \gamma_1 + \gamma_2) + (\beta_1 + \delta_1 + \delta_2)\,x_{t_1}
      & \gamma_2 + \delta_2x_{t_1} \\
\bottomrule
\end{tabular}
\end{center}
При вычитании всё общее у двух прямых сокращается, и остаётся добавка нового
промежутка $\gamma_j + \delta_jx$ в точке сдвига. Здесь $\gamma_j$ — расстояние
между прямыми при $x = 0$, а $\delta_j$ — на сколько оно меняется при росте $x$
на единицу. Поэтому скачок \emph{не} равен $\gamma_j$: даже при $\gamma_j = 0$
одна лишь смена наклона даёт скачок $\delta_jx_{t_j}$.

\begin{center}
\begin{tikzpicture}[>=stealth, line join=round, x=0.95cm, y=0.95cm]
  % x_t = t; t0 = 3, t1 = 6; наклоны 0,4 / 0,6 / 0,1; скачки -0,6 и +0,6
  \draw[->] (0,0) -- (9.8,0) node[anchor=north] {$t$};
  \draw[->] (0,0) -- (0,5.6) node[anchor=east]  {$y$};
  \draw[dotted, black!55] (3,0) -- (3,5.2);
  \draw[dotted, black!55] (6,0) -- (6,5.2);
  \node[anchor=north] at (3,0) {$t_0$};
  \node[anchor=north] at (6,0) {$t_1$};

  % продолжения старых прямых — от них отсчитывается скачок
  \draw[dashed, accent!55] (3,2.2) -- (4.2,2.68);
  \draw[dashed, solbar!55] (6,3.4) -- (7.2,4.12);

  % три отрезка
  \draw[very thick, accent]  (0,1.0) -- (3,2.2);
  \draw[very thick, solbar]  (3,1.6) -- (6,3.4);
  \draw[very thick, notebar] (6,4.0) -- (9,4.3);
  \fill[accent] (0,1.0) circle (1.6pt);
  \node[anchor=east, accent] at (0,1.0) {$\beta_0$};

  % наклоны
  \node[accent,  anchor=south east] at (2.6,2.15) {наклон $\beta_1$};
  \node[solbar,  anchor=north west] at (4.4,2.35) {наклон $\beta_1 + \delta_1$};
  \node[notebar, anchor=south]      at (7.6,4.3)  {наклон $\beta_1 + \delta_1 + \delta_2$};

  % скачки
  \draw[<->, thick] (3.12,2.2) -- (3.12,1.6);
  \node[anchor=north, font=\small] at (3.12,1.55) {$\gamma_1 + \delta_1t_0$};
  \draw[<->, thick] (6.12,3.4) -- (6.12,4.0);
  \node[anchor=east, font=\small] at (5.95,4.05) {$\gamma_2 + \delta_2t_1$};
\end{tikzpicture}

\smallskip
{\small\itshape На рисунке $x_t = t$ (линейный тренд). Цвет отрезка — промежуток;
пунктир — продолжение прямой предыдущего промежутка, стрелка — скачок в момент
сдвига.\par}
\end{center}

\medskip
\textbf{3) Проверка гипотез} — по \eqref{eq:restr}, в модели $k = 6$
параметров:
\begin{center}
\begin{tabular}{lll}
\toprule
Что проверяем & $H_0$ & $q$ \\
\midrule
нет сдвига в момент $t_0$ & $\gamma_1 = \delta_1 = 0$ & 2 \\
нет сдвига в момент $t_1$ & $\gamma_2 = \delta_2 = 0$ & 2 \\
нет ни одного сдвига      & $\gamma_1 = \delta_1 = \gamma_2 = \delta_2 = 0$ & 4 \\
\bottomrule
\end{tabular}
\end{center}
Последняя гипотеза — тест Чоу \eqref{eq:chow} для $p = 3$ частей: модель с
фиктивными переменными подгоняет три промежутка независимо, так что
$\ESS_{\mathrm{UR}} = \ESS_1 + \ESS_2 + \ESS_3$, а $\ESS_{\mathrm{R}}$ —
остаточная сумма одной регрессии по всей выборке:
\[
  F = \frac{\bigl(\ESS_{\mathrm{R}} - \ESS_{\mathrm{UR}}\bigr)/q}
           {\ESS_{\mathrm{UR}}/(n-k)}
    = \frac{\bigl(\ESS_{\mathrm{R}} - (\ESS_1 + \ESS_2 + \ESS_3)\bigr)/4}
           {(\ESS_1 + \ESS_2 + \ESS_3)/(n - 6)} \sim F(4,\ n-6).
\]
\end{solution}

% -------------------------------------------------------------------
\ifsolutions\clearpage\fi
\problem{10.4 (Бородич)}

В следующей таблице приведены данные по реальному валовому национальному
продукту ($\mathit{GNP}$), реальному объёму потребления ($\mathit{CONS}$) и
объёму инвестиций ($\mathit{INV}$) для некоторой вымышленной страны.

\begin{center}
\small
\setlength{\tabcolsep}{4pt}
\begin{tabular}{l|*{10}{r}}
$\mathit{GNP}$  & 240 & 248 & 261 & 274 & 273 & 269 & 283 & 296 & 312 & 319 \\
\hline
$\mathit{CONS}$ & 149 & 154 & 162 & 169 & 167 & 171 & 180 & 188 & 196 & 200 \\
\hline
$\mathit{INV}$  & 38,2 & 41,9 & 46,5 & 52,1 & 48,1 & 38,3 & 45,4 & 52,1 & 56,8 & 57,5 \\
\multicolumn{11}{c}{} \\[-6pt]
$\mathit{GNP}$  & 318 & 325 & 317 & 327 & 350 & 361 & 372 & 385 & 402 & 412 \\
\hline
$\mathit{CONS}$ & 200 & 202 & 205 & 215 & 225 & 235 & 245 & 252 & 261 & 266 \\
\hline
$\mathit{INV}$  & 50,9 & 54,5 & 44,7 & 50,4 & 65,8 & 63,7 & 64,0 & 76,4 & 71,6 & 71,8 \\
\end{tabular}
\end{center}

\begin{enumerate}[label=\asbuk*), leftmargin=2.2em, itemsep=3pt]
  \item Постройте уравнение регрессии
        $\mathit{INV} = b_0 + b_1\mathit{GNP} + b_2\mathit{CONS} + e$.
  \item Оцените качество построенного уравнения.
  \item Можно ли было ожидать при построении данного уравнения наличия
        мультиколлинеарности? Ответ поясните.
  \item Имеет ли место мультиколлинеарность для построенного вами уравнения?
        Как вы это определили?
  \item Постройте уравнения регрессии $\mathit{INV}$ на $\mathit{GNP}$ и
        $\mathit{INV}$ на $\mathit{CONS}$. Какие выводы можно сделать по
        построенным моделям?
  \item Постройте уравнение регрессии $\mathit{CONS}$ на $\mathit{GNP}$. Что
        обнаруживает построенная модель?
  \item Как можно решить проблему мультиколлинеарности для первоначальной
        модели?
\end{enumerate}

\begin{excel}
Введите данные тремя столбцами $\mathit{GNP}$, $\mathit{CONS}$, $\mathit{INV}$
(20 строк). Регрессию с двумя регрессорами даёт \ui{Данные} $\to$
\ui{Анализ данных} $\to$ \ui{Регрессия}: входной интервал $Y$ — столбец
$\mathit{INV}$, входной интервал $X$ — \emph{оба} столбца $\mathit{GNP}$ и
$\mathit{CONS}$ (они должны стоять рядом). Матрицу корреляций строит
\ui{Анализ данных} $\to$ \ui{Корреляция}, отдельный коэффициент —
\xl{=КОРРЕЛ(\dots)}. Коэффициент детерминации парной регрессии — \xl{=КВПИРСОН(y; x)},
отсюда VIF~= \xl{=1/(1-КВПИРСОН(CONS; GNP))}.
\end{excel}

\begin{solution}
Здесь $n = 20$. Критические значения: $t_{0{,}025}(17) = 2{,}110$,
$t_{0{,}025}(18) = 2{,}101$, $F_{0{,}05}(2, 17) = 3{,}592$,
$F_{0{,}05}(1, 18) = 4{,}414$.

\medskip
\textbf{а) Регрессия на оба регрессора} ($k = 3$, $n - k = 17$):
\begin{center}
\begin{tabular}{lrrrr}
\toprule
Переменная & Коэффициент & Ст. ошибка & $t$-статистика & $P$-значение \\
\midrule
const           & $-22{,}593$ & 9,524 & $-2{,}372$ & 0,0297 \\
$\mathit{GNP}$  & $\phantom{-}0{,}602$ & 0,231 & $\phantom{-}2{,}602$ & 0,0186 \\
$\mathit{CONS}$ & $-0{,}563$ & 0,326 & $-1{,}725$ & 0,1026 \\
\midrule
\multicolumn{5}{l}{$R^2 = 0{,}8681$;\quad $R^2_{\text{adj}} = 0{,}8525$;\quad
  $\ESS = 316{,}65$;\quad $\TSS = 2400{,}15$} \\
\bottomrule
\end{tabular}
\end{center}
\[
  \widehat{\mathit{INV}} = -22{,}59 + 0{,}602\,\mathit{GNP} - 0{,}563\,\mathit{CONS}.
\]

\medskip
\textbf{б) Качество уравнения.} Регрессия в целом значима: по \eqref{eq:tf}
\[
  F = \frac{R^2}{1-R^2}\cdot\frac{n-k}{k-1}
    = \frac{0{,}8681}{0{,}1319}\cdot\frac{17}{2}
    = 55{,}93 > 3{,}592 = F_{0{,}05}(2, 17),
  \qquad P = 3{,}3\cdot 10^{-8}.
\]
Уравнение объясняет 87\,\% разброса инвестиций. Но по отдельности картина
хуже: коэффициент при $\mathit{GNP}$ значим ($|t| = 2{,}60 > 2{,}110$), а при
$\mathit{CONS}$ — нет ($|t| = 1{,}73 < 2{,}110$, $P = 0{,}10$). Знак при
$\mathit{CONS}$ отрицательный. Формально это можно истолковать: при
\emph{неизменном} ВНП рост потребления означает меньшие сбережения, а значит,
и меньшие инвестиции. Но такой «эксперимент» — $\mathit{CONS}$ меняется, а
$\mathit{GNP}$ стоит на месте — в данных почти не встречается, поэтому оценка
опирается на очень малую вариацию и неточна.

\medskip
\textbf{в) Можно ли было ожидать мультиколлинеарности?} Да. Потребление —
основная часть ВНП ($\mathit{GNP} = \mathit{CONS} + \mathit{INV} + \dots$), а
по кейнсианской функции потребления $\mathit{CONS}$ линейно зависит от дохода.
Кроме того, обе переменные — растущие во времени макроэкономические ряды. Всё
это заранее говорит о тесной линейной связи двух регрессоров.

\medskip
\textbf{г) Есть ли мультиколлинеарность?} Да, по всем признакам
п.~\ref{sec:mc}.
\begin{itemize}[leftmargin=2.2em, itemsep=2pt]
  \item Корреляционная матрица:
        \[
          \begin{array}{l|ccc}
                          & \mathit{INV} & \mathit{GNP} & \mathit{CONS} \\ \hline
            \mathit{INV}  & 1       & 0{,}919 & 0{,}903 \\
            \mathit{GNP}  & 0{,}919 & 1       & 0{,}996 \\
            \mathit{CONS} & 0{,}903 & 0{,}996 & 1
          \end{array}
        \]
        Корреляция регрессоров $r(\mathit{GNP}, \mathit{CONS}) = 0{,}996$ выше,
        чем корреляция каждого из них с $\mathit{INV}$.
  \item VIF. Регрессоров два, поэтому вспомогательная регрессия одна — это
        регрессия $\mathit{CONS}$ на $\mathit{GNP}$ из п.~е), её
        $R_j^2 = 0{,}9928$. По \eqref{eq:vif}
        \[
          \mathrm{VIF} = \frac{1}{1 - R_j^2} = \frac{1}{1 - 0{,}9928} = 139{,}7 \gg 10 .
        \]
  \item $F$-тест вспомогательной регрессии \eqref{eq:fj} при $m = 2$
        объясняющих переменных:
        \[
          F_j = \frac{R_j^2}{1 - R_j^2}\cdot\frac{n-m}{m-1}
              = \frac{0{,}9928}{0{,}0072}\cdot\frac{18}{1}
              = 2496 \gg 4{,}414 = F_{0{,}05}(1, 18).
        \]
  \item Типичный симптом: высокий $R^2$ и значимый $F$ при незначимом
        отдельном коэффициенте. Стандартная ошибка при $\mathit{GNP}$ в
        регрессии на оба регрессора (0,231) в 11 раз больше, чем в парной
        регрессии из п.~д) (0,021), — это почти ровно
        $\sqrt{\mathrm{VIF}} = 11{,}8$.
\end{itemize}

\medskip
\textbf{д) Парные регрессии} ($k = 2$, $n - k = 18$):
\begin{center}
\begin{tabular}{lrrrr}
\toprule
Переменная & Коэффициент & Ст. ошибка & $t$-статистика & $P$-значение \\
\midrule
\multicolumn{5}{l}{\emph{$\mathit{INV}$ на $\mathit{GNP}$}} \\
const          & $-10{,}242$ & 6,619 & $-1{,}548$ & 0,1391 \\
$\mathit{GNP}$ & $\phantom{-}0{,}204$ & 0,021 & $9{,}905$ & 0,0000 \\
\multicolumn{5}{l}{$R^2 = 0{,}8450$;\quad $R^2_{\text{adj}} = 0{,}8364$;\quad $\ESS = 372{,}11$} \\
\midrule
\multicolumn{5}{l}{\emph{$\mathit{INV}$ на $\mathit{CONS}$}} \\
const           & $-2{,}667$ & 6,508 & $-0{,}410$ & 0,6868 \\
$\mathit{CONS}$ & $\phantom{-}0{,}283$ & 0,032 & $8{,}920$ & 0,0000 \\
\multicolumn{5}{l}{$R^2 = 0{,}8155$;\quad $R^2_{\text{adj}} = 0{,}8053$;\quad $\ESS = 442{,}78$} \\
\bottomrule
\end{tabular}
\end{center}
Выводы:
\begin{itemize}[leftmargin=2.2em, itemsep=2pt]
  \item По отдельности каждый регрессор высоко значим ($t \approx 9$--$10$), и
        знак у обоих положительный, как и подсказывает экономика.
  \item $R^2$ почти не отличается от $R^2$ регрессии на оба регрессора
        ($0{,}845$ против $0{,}868$): второй регрессор несёт мало новой
        информации — почти всё, что он знает об $\mathit{INV}$, уже содержится
        в первом.
  \item Добавление $\mathit{CONS}$ к регрессии на $\mathit{GNP}$ — это одно
        ограничение $b_2 = 0$. По \eqref{eq:restr} с
        $\ESS_{\mathrm{R}} = 372{,}11$ (парная) и $\ESS_{\mathrm{UR}} = 316{,}65$
        (из п.~а):
        \[
          F = \frac{\bigl(\ESS_{\mathrm{R}} - \ESS_{\mathrm{UR}}\bigr)/q}
                   {\ESS_{\mathrm{UR}}/(n-k)}
            = \frac{(372{,}11 - 316{,}65)/1}{316{,}65/17}
            = 2{,}98 < 4{,}451 = F_{0{,}05}(1, 17),
        \]
        и это в точности $t^2 = 1{,}725^2$ для $\mathit{CONS}$ из п.~а).
        Добавлять $\mathit{CONS}$ смысла нет.
  \item Почему наклон при $\mathit{GNP}$ в парной регрессии (0,204) так
        отличается от 0,602 в п.~а)? По тождеству \eqref{eq:short} с
        $\widehat{\delta} = 0{,}706$ из п.~е):
        \[
          \widetilde{b}_1 = \widehat{b}_1 + \widehat{b}_2\,\widehat{\delta}
                          = 0{,}602 + (-0{,}563)\cdot 0{,}706 = 0{,}204 .
        \]
        Парная регрессия показывает \emph{суммарный} эффект роста ВНП вместе с
        сопутствующим ростом потребления, а 0,602 — эффект при неизменном
        потреблении.
\end{itemize}

\medskip
\textbf{е) Регрессия $\mathit{CONS}$ на $\mathit{GNP}$:}
\[
  \widehat{\mathit{CONS}} = \underset{(4{,}538)}{-21{,}94}
                          + \underset{(0{,}014)}{0{,}706}\,\mathit{GNP},
  \qquad t = 49{,}96,\quad R^2 = 0{,}9928 .
\]
Модель обнаруживает почти точную линейную зависимость между регрессорами
исходного уравнения: $\mathit{GNP}$ объясняет 99,3\,\% разброса
$\mathit{CONS}$. Экономически это функция потребления с предельной склонностью
к потреблению 0,706. Для задачи о мультиколлинеарности это и есть
вспомогательная регрессия, по которой в п.~г) посчитан VIF.

\begin{center}
\begin{tikzpicture}[>=stealth, x=0.8cm, y=0.8cm]
  % оси: x = (GNP - 230)/25, y = (CONS - 140)/20
  \draw[->] (0,0) -- (8.0,0) node[anchor=north west] {$\mathit{GNP}$};
  \draw[->] (0,0) -- (0,7.0) node[anchor=east] {$\mathit{CONS}$};
  \foreach \g/\x in {250/0.8, 300/2.8, 350/4.8, 400/6.8}
    \draw (\x,0) -- (\x,-0.1) node[anchor=north, font=\small] {\g};
  \foreach \c/\y in {150/0.5, 200/3.0, 250/5.5}
    \draw (0,\y) -- (-0.1,\y) node[anchor=east, font=\small] {\c};

  % линия регрессии CONS = -21,94 + 0,706 GNP
  \draw[thick, accent] (0.2,0.20) -- (7.4,6.56)
        node[anchor=west, font=\small] {$\widehat{\mathit{CONS}}$};
  \foreach \p in {(0.40,0.45),(0.72,0.70),(1.24,1.10),(1.76,1.45),(1.72,1.35),
                  (1.56,1.55),(2.12,2.00),(2.64,2.40),(3.28,2.80),(3.56,3.00),
                  (3.52,3.00),(3.80,3.10),(3.48,3.25),(3.88,3.75),(4.80,4.25),
                  (5.24,4.75),(5.68,5.25),(6.20,5.60),(6.88,6.05),(7.28,6.30)}
    \fill[solbar] \p circle (2pt);
\end{tikzpicture}

\smallskip
{\small\itshape Все 20 точек лежат почти на одной прямой ($R^2 = 0{,}993$):
у данных нет наблюдений, где $\mathit{CONS}$ заметно меняется при том же
$\mathit{GNP}$. Именно такие наблюдения нужны, чтобы разделить влияние двух
регрессоров.\par}
\end{center}

\medskip
\textbf{ж) Как устранить мультиколлинеарность.} Проблема в том, что
$\mathit{GNP}$ и $\mathit{CONS}$ почти пропорциональны ($r = 0{,}996$), и модель
не может разделить их влияние. Рассмотрим два способа из п.~\ref{sec:cure}.

\medskip
\emph{Способ 1: исключить $\mathit{CONS}$.} Остаётся регрессия $\mathit{INV}$
на $\mathit{GNP}$ из п.~д): коэффициент значим, знак правильный,
$R^2_{\text{adj}}$ снизился лишь с 0,853 до 0,836. Минус в том, что её
коэффициент 0,204 — суммарный эффект роста ВНП вместе с сопутствующим ростом
потребления (п.~д), а не эффект ВНП при прочих равных.

\medskip
\emph{Способ 2: заменить два регрессора одним осмысленным.}

\smallskip
\textbf{Шаг 1. Наблюдение.} В п.~а) $\widehat{b}_1 = 0{,}602$ и
$\widehat{b}_2 = -0{,}563$ — почти одинаковы по модулю и противоположны по
знаку. Если бы было в точности $b_2 = -b_1$, то
\[
  b_1\,\mathit{GNP} + b_2\,\mathit{CONS}
  = b_1\,(\mathit{GNP} - \mathit{CONS}) = b_1S,
  \qquad S = \mathit{GNP} - \mathit{CONS},
\]
то есть инвестиции зависели бы не от $\mathit{GNP}$ и $\mathit{CONS}$ по
отдельности, а только от их разности $S$ — сбережений. Это и есть одна
переменная вместо двух. Но сначала надо проверить, что данные не
противоречат ограничению $H_0\colon b_1 + b_2 = 0$.

\smallskip
\textbf{Шаг 2. Проверка ограничения.} Перепишем исходную модель, ничего в ней
не меняя, — прибавим и вычтем $b_2\,\mathit{GNP}$:
\[
  b_1\,\mathit{GNP} + b_2\,\mathit{CONS}
  = (b_1 + b_2)\,\mathit{GNP} - b_2\,(\mathit{GNP} - \mathit{CONS})
  = (b_1 + b_2)\,\mathit{GNP} - b_2S .
\]
Значит, регрессия $\mathit{INV}$ на $\mathit{GNP}$ и $S$ — та же самая модель:
подгонка и $\ESS = 316{,}65$ совпадают с п.~а). Зато коэффициент при
$\mathit{GNP}$ в ней равен ровно $b_1 + b_2$, и гипотеза $H_0\colon b_1 + b_2 = 0$
проверяется обычным $t$-тестом этого коэффициента:
\[
  \widehat{b_1 + b_2} = 0{,}039,
  \quad s = 0{,}098,
  \quad t = \frac{0{,}039}{0{,}098} = 0{,}40 < 2{,}110 = t_{0{,}025}(17),
  \quad P = 0{,}70 .
\]
То же даёт $F$-тест \eqref{eq:restr} с $q = 1$: модель с ограничением —
регрессия на одно $S$ (шаг~3), $\ESS_{\mathrm{R}} = 319{,}59$, без
ограничения — модель из п.~а), $\ESS_{\mathrm{UR}} = 316{,}65$:
\[
  F = \frac{\bigl(\ESS_{\mathrm{R}} - \ESS_{\mathrm{UR}}\bigr)/q}
           {\ESS_{\mathrm{UR}}/(n-k)}
    = \frac{(319{,}59 - 316{,}65)/1}{316{,}65/17}
    = 0{,}16 = t^2 < 4{,}451 = F_{0{,}05}(1, 17).
\]
Ограничение не отвергается: данные согласуются с тем, что важна только
разность $\mathit{GNP} - \mathit{CONS}$.

\smallskip
\textbf{Шаг 3. Новая модель.} Подставляем ограничение $b_2 = -b_1$:
$\mathit{INV} = b_0 + b_1S + e$. Оценка:
\[
  \widehat{\mathit{INV}} = \underset{(7{,}397)}{-24{,}89}
                         + \underset{(0{,}064)}{0{,}690}\,S,
  \qquad t = 10{,}83,\quad R^2 = 0{,}8668,\quad R^2_{\text{adj}} = 0{,}8594 .
\]
Регрессор один, так что мультиколлинеарности нет, и коэффициент оценён точно.
Сравнение всех моделей задачи:
\begin{center}
\begin{tabular}{lcc}
\toprule
Модель & $R^2_{\text{adj}}$ & Проблема \\
\midrule
$\mathit{INV}$ на $\mathit{GNP}$ и $\mathit{CONS}$ & 0,8525 & мультиколлинеарность \\
$\mathit{INV}$ на $\mathit{GNP}$                  & 0,8364 & коэффициент — суммарный эффект \\
$\mathit{INV}$ на $\mathit{CONS}$                 & 0,8053 & то же \\
$\mathit{INV}$ на $S$                             & 0,8594 & — \\
\bottomrule
\end{tabular}
\end{center}
Модель с $S$ лучшая по $R^2_{\text{adj}}$.

\begin{remark}
Сбережения — экономически осмысленная переменная: инвестиции финансируются
из сбережений. Но $S = \mathit{GNP} - \mathit{CONS}$ содержит $\mathit{INV}$
как слагаемое ($S = \mathit{INV}$ + госрасходы + чистый экспорт), поэтому
связь $\mathit{INV}$ с $S$ отчасти следует из того, как устроены сами
переменные, и высокий $R^2$ здесь не открытие.
\end{remark}

\medskip
\emph{Другие способы:} увеличить выборку; перейти к приростам
($\Delta\mathit{INV}$ на $\Delta\mathit{GNP}$ и $\Delta\mathit{CONS}$) —
приросты связаны слабее, чем уровни растущих рядов; перейти к относительным
величинам ($\mathit{INV}/\mathit{GNP}$ на $\mathit{CONS}/\mathit{GNP}$).
\end{solution}

% -------------------------------------------------------------------
\ifsolutions\clearpage\fi
\problem{4.17}

В таблице представлены выпуск $Q$, трудозатраты $L$ и капиталовложения $K$
15 фирм некоторой отрасли.

\begin{center}
\begin{tabular}{cccc @{\qquad} cccc}
\toprule
Фирма & $Q$ & $L$ & $K$ & Фирма & $Q$ & $L$ & $K$ \\
\midrule
1 & 2350 & 2334 & 1570 &  9 & 2550 & 2446 & 1880 \\
2 & 2470 & 2425 & 1850 & 10 & 2450 & 2403 & 1790 \\
3 & 2110 & 2230 & 1150 & 11 & 2290 & 2301 & 1480 \\
4 & 2560 & 2463 & 1940 & 12 & 2160 & 2253 & 1240 \\
5 & 2650 & 2565 & 2450 & 13 & 2400 & 2367 & 1660 \\
6 & 2240 & 2278 & 1340 & 14 & 2490 & 2430 & 1850 \\
7 & 2430 & 2380 & 1700 & 15 & 2590 & 2470 & 2000 \\
8 & 2530 & 2437 & 1860 &    &      &      &      \\
\bottomrule
\end{tabular}
\end{center}

\begin{enumerate}[label=\asbuk*), leftmargin=2.2em, itemsep=3pt]
  \item Оцените по этим данным производственную функцию Кобба--Дугласа
        $Q = \alpha L^{\beta_1}K^{\beta_2}$, вычислите коэффициент детерминации,
        скорректированный коэффициент детерминации и выборочный коэффициент
        корреляции между $\ln L$ и $\ln K$.
  \item Проведите регрессию $\ln Q$ только на $\ln K$.
\end{enumerate}
Как можно проинтерпретировать полученные результаты?

\begin{excel}
Добавьте три столбца логарифмов: \xl{=LN(B2)} и т.\,д. Регрессия
$\ln Q$ на $\ln L$ и $\ln K$ — \ui{Анализ данных} $\to$ \ui{Регрессия} с
входным интервалом $X$ из двух соседних столбцов $\ln L$, $\ln K$. Корреляция —
\xl{=КОРРЕЛ(lnL; lnK)}.
\end{excel}

\begin{solution}
Здесь $n = 15$. Критические значения: $t_{0{,}025}(12) = 2{,}179$,
$t_{0{,}025}(13) = 2{,}160$, $F_{0{,}05}(2, 12) = 3{,}885$.

\medskip
\textbf{а) Функция Кобба--Дугласа.} Функция $Q = \alpha L^{\beta_1}K^{\beta_2}$
нелинейна по параметрам, поэтому сначала её логарифмируем. Логарифм
произведения — сумма логарифмов, а логарифм степени — показатель, умноженный на
логарифм основания:
\[
  \ln Q = \ln\bigl(\alpha L^{\beta_1}K^{\beta_2}\bigr)
        = \ln\alpha + \ln L^{\beta_1} + \ln K^{\beta_2}
        = \ln\alpha + \beta_1\ln L + \beta_2\ln K .
\]
Добавив случайную ошибку, получаем линейную регрессию \eqref{eq:cd}
\[
  \ln Q_i = \beta_0 + \beta_1\ln L_i + \beta_2\ln K_i + \varepsilon_i,
  \qquad \beta_0 = \ln\alpha .
\]
Зависимая переменная — $\ln Q$, регрессоры — $\ln L$ и $\ln K$. Модель линейна
по параметрам $\beta_0$, $\beta_1$, $\beta_2$, и её оценивает обычный МНК
($k = 3$, $n - k = 12$). Сам параметр $\alpha$ восстанавливается как
$\widehat{\alpha} = e^{\widehat{\beta}_0}$.
\begin{center}
\begin{tabular}{llrrrr}
\toprule
Регрессор & Параметр & Оценка & Ст. ошибка & $t$-статистика & $P$-значение \\
\midrule
константа & $\beta_0 = \ln\alpha$ & 0,501 & 4,480 & 0,112 & 0,9129 \\
$\ln L$   & $\beta_1$            & 0,758 & 0,707 & 1,071 & 0,3052 \\
$\ln K$   & $\beta_2$            & 0,188 & 0,139 & 1,356 & 0,2001 \\
\midrule
\multicolumn{6}{l}{$R^2 = 0{,}9689$} \\
\bottomrule
\end{tabular}
\end{center}
\[
  \widehat{Q} = e^{0{,}501}L^{0{,}758}K^{0{,}188} = 1{,}65\,L^{0{,}758}K^{0{,}188}.
\]
Скорректированный коэффициент детерминации по \eqref{eq:r2adj}:
\[
  R^2_{\text{adj}} = 1 - \frac{n-1}{n-k}\bigl(1 - R^2\bigr)
                   = 1 - \frac{14}{12}\cdot(1 - 0{,}9689) = 0{,}9637 .
\]
То же через суммы квадратов ($\ESS = 0{,}001981$, $\TSS = 0{,}06366$) и через
третью запись \eqref{eq:r2adj}:
\[
  R^2_{\text{adj}} = 1 - \frac{\ESS/(n-k)}{\TSS/(n-1)}
                   = 1 - \frac{0{,}001981/12}{0{,}06366/14} = 0{,}9637,
\]
\[
  R^2_{\text{adj}} = R^2 - \frac{k-1}{n-k}\bigl(1 - R^2\bigr)
                   = 0{,}9689 - \frac{2}{12}\cdot 0{,}0311 = 0{,}9637 .
\]
Выборочный коэффициент корреляции $r(\ln L, \ln K) = 0{,}992$.

\medskip
\textbf{Интерпретация.} Классическая картина мультиколлинеарности:
\begin{itemize}[leftmargin=2.2em, itemsep=2pt]
  \item регрессия в целом значима и объясняет почти 97\,\% разброса
        $\ln Q$: $F = \frac{0{,}9689}{0{,}0311}\cdot\frac{12}{2} = 186{,}8
        \gg 3{,}885$;
  \item при этом \emph{ни одна} эластичность не значима: $|t| = 1{,}07$ и
        $1{,}36$ меньше $2{,}179$;
  \item причина — почти линейная связь $\ln L$ и $\ln K$: фирмы с большим
        капиталом нанимают и больше труда. По \eqref{eq:vif}
        \[
          \mathrm{VIF} = \frac{1}{1 - r^2} = \frac{1}{1 - 0{,}992^2} = 63{,}7 ,
        \]
        то есть дисперсии оценок в 64 раза больше, чем при несвязанных
        факторах.
\end{itemize}
Данные хорошо описывают выпуск в целом, но не позволяют разделить вклады
труда и капитала.

\medskip
\textbf{б) Регрессия $\ln Q$ только на $\ln K$} ($n - k = 13$):
\begin{center}
\begin{tabular}{lrrrr}
\toprule
Переменная & Коэффициент & Ст. ошибка & $t$-статистика & $P$-значение \\
\midrule
const   & 5,297 & 0,130 & 40,778 & 0,0000 \\
$\ln K$ & 0,335 & 0,017 & 19,192 & 0,0000 \\
\midrule
\multicolumn{5}{l}{$R^2 = 0{,}9659$} \\
\bottomrule
\end{tabular}
\end{center}
Коэффициент стал высоко значимым, а $R^2$ почти не изменился
($0{,}9659$ против $0{,}9689$): $\ln L$ добавлял к $\ln K$ очень мало
информации.

\medskip
\textbf{Как это понимать.}
\begin{itemize}[leftmargin=2.2em, itemsep=2pt]
  \item Для \emph{прогноза} выпуска короткой регрессии достаточно: она почти
        так же точна и гораздо проще.
  \item Но 0,335 — \emph{не} эластичность выпуска по капиталу. По
        \eqref{eq:short}, где $\widehat{\delta} = 0{,}195$ — наклон регрессии
        $\ln L$ на $\ln K$,
        \[
          0{,}188 + 0{,}758\cdot 0{,}195 = 0{,}335 :
        \]
        коэффициент включает и эффект труда, который растёт вместе с
        капиталом. Чтобы оценить эластичности по труду и капиталу по
        отдельности, нужна регрессия на оба фактора, а она из-за
        мультиколлинеарности даёт их с большими ошибками. Один из выходов —
        следующая задача.
\end{itemize}
\end{solution}

% -------------------------------------------------------------------
\problem{4.18}

Можно ли преодолеть проблему мультиколлинеарности, возникающую в
задаче~4.17, если известно, что производственная функция обладает постоянной
отдачей на масштаб ($\beta_1 + \beta_2 = 1$)?

\begin{excel}
В той же книге, что и для задачи~4.17, добавьте столбцы $\ln(Q/K)$ и
$\ln(L/K)$ и оцените парную регрессию первого на второй. Остаточные суммы для
$F$-теста — строка \ui{Остаток}, столбец \ui{SS} в обоих отчётах.
\end{excel}

\begin{solution}
Да. Ограничение — это внешняя информация, которой нет в данных, и она
заменяет недостающую вариацию регрессоров (п.~\ref{sec:cure}, способ 3).

\medskip
\textbf{1) Проверка ограничения.} Сначала убедимся, что данные не
противоречат гипотезе $H_0\colon \beta_1 + \beta_2 = 1$ ($q = 1$). Точечная
оценка суммы из задачи~4.17: $0{,}758 + 0{,}188 = 0{,}946$ — близко к единице.

Модель с ограничением получаем подстановкой $\beta_2 = 1 - \beta_1$ в
\eqref{eq:cd}:
\[
  \ln Q = \ln\alpha + \beta_1\ln L + (1 - \beta_1)\ln K + \varepsilon
  \quad\Longrightarrow\quad
  \ln\frac{Q}{K} = \ln\alpha + \beta_1\ln\frac{L}{K} + \varepsilon .
\]
Выпуск на единицу капитала объясняется трудом на единицу капитала. Оценка:
\begin{center}
\begin{tabular}{lrrrr}
\toprule
Переменная & Коэффициент & Ст. ошибка & $t$-статистика & $P$-значение \\
\midrule
const       & 0,073 & 0,008 & \phantom{0}9,261 & 0,0000 \\
$\ln(L/K)$  & 0,825 & 0,021 & 39,812 & 0,0000 \\
\midrule
\multicolumn{5}{l}{$R^2 = 0{,}9919$;\quad
  $\ESS_{\mathrm{R}} = 1{,}9826\cdot 10^{-3}$} \\
\bottomrule
\end{tabular}
\end{center}
Модель без ограничений — регрессия из задачи~4.17, а),
$\ESS_{\mathrm{UR}} = 1{,}9811\cdot 10^{-3}$, $n - k = 15 - 3 = 12$. По
\eqref{eq:restr}
\[
  F = \frac{\bigl(\ESS_{\mathrm{R}} - \ESS_{\mathrm{UR}}\bigr)/q}
           {\ESS_{\mathrm{UR}}/(n-k)}
    = \frac{\bigl(\ESS_{\mathrm{R}} - \ESS_{\mathrm{UR}}\bigr)/1}
           {\ESS_{\mathrm{UR}}/12},
\]
\[
  F = \frac{(1{,}9826 - 1{,}9811)\cdot 10^{-3}}{1{,}9811\cdot 10^{-3}/12}
    = 0{,}009 < 4{,}747 = F_{0{,}05}(1, 12),
  \qquad P = 0{,}926 .
\]
(Разность остаточных сумм в числителе точнее равна $1{,}51\cdot 10^{-6}$: при
четырёх знаках после запятой она почти теряется.) Гипотеза о постоянной отдаче
на масштаб уверенно не отвергается.

\begin{remark}
Коэффициент детерминации модели с ограничением ($0{,}9919$) \emph{больше}, чем
у модели без ограничений ($0{,}9689$), хотя ограничение не может улучшить
подгонку. Противоречия нет: у моделей разные зависимые переменные, $\ln(Q/K)$
и $\ln Q$, поэтому и $\TSS$ разные ($0{,}2437$ против $0{,}0637$), и их $R^2$
несравнимы. Сравнивать надо остаточные суммы, как в формуле выше.
\end{remark}

\medskip
\textbf{2) Результат.} Оценки параметров:
\[
  \widehat{\beta}_1 = 0{,}825,
  \qquad
  \widehat{\beta}_2 = 1 - \widehat{\beta}_1 = 0{,}175,
  \qquad
  \widehat{\alpha} = e^{0{,}073} = 1{,}075,
\]
\[
  \widehat{Q} = 1{,}075\,L^{0{,}825}K^{0{,}175}.
\]
Стандартная ошибка эластичности по труду упала с 0,707 до 0,021, в 34 раза, а
$t$-статистика выросла с 1,07 до 39,8. У $\widehat{\beta}_2$ та же стандартная
ошибка 0,021, так как $\widehat{\beta}_2 = 1 - \widehat{\beta}_1$. В модели
остался один регрессор $\ln(L/K)$, так что мультиколлинеарности нет.

\medskip
\textbf{Вывод.} Введение ограничения $\beta_1 + \beta_2 = 1$ решило проблему
мультиколлинеарности: обе эластичности оценены точно. Цена — предположение о
постоянной отдаче на масштаб. Если бы оно было неверным, оценки были бы
смещены, поэтому ограничение и нужно было сначала проверить.
\end{solution}

% -------------------------------------------------------------------
\ifsolutions\clearpage\fi
\problem{4.1.1 (задачник ВШЭ)}

Дана матрица
\[
  X = \begin{pmatrix}
        1 & 10 & 3 \\
        1 &  6 & 1 \\
        1 & 12 & a \\
        1 &  8 & 2
      \end{pmatrix}.
\]
Каким должно быть число $a$, чтобы между столбцами матрицы была строгая
мультиколлинеарность? Допустим, $a = 2$. Рассчитайте значения $\mathrm{VIF}$
для оценок коэффициентов регрессии $y = X\beta + \varepsilon$, где $y$ —
любая величина.

\begin{solution}
Обозначим столбцы $\mathbf{1}$, $\mathbf{x}_2$, $\mathbf{x}_3$.

\medskip
\textbf{1) Строгая мультиколлинеарность} — линейная зависимость столбцов
(п.~\ref{sec:mc}). Столбцы $\mathbf{1}$ и $\mathbf{x}_2$ не пропорциональны,
поэтому зависимость возможна только в виде
$\mathbf{x}_3 = c_1\mathbf{1} + c_2\mathbf{x}_2$. Коэффициенты находим по
строкам, где $\mathbf{x}_3$ известен. Из первых двух строк
\[
  \begin{cases}
    c_1 + 10c_2 = 3, \\
    c_1 + \phantom{0}6c_2 = 1
  \end{cases}
  \quad\Longrightarrow\quad
  c_2 = \tfrac12,\quad c_1 = -2 .
\]
Проверка по четвёртой строке: $-2 + \tfrac12\cdot 8 = 2$ — верно. Значит,
при $x_3 = -2 + \tfrac12x_2$ третья строка даёт
\[
  a = -2 + \tfrac12\cdot 12 = 4 .
\]
При $a = 4$ ранг $X$ равен 2, $\det(X^{\top}X) = 0$, и МНК-оценка не
существует. При любом другом $a$ столбцы независимы.

\medskip
\textbf{2) VIF при $a = 2$.} VIF зависит только от матрицы регрессоров, а не
от $y$, — поэтому $y$ в условии и может быть любым. Регрессоров (кроме
константы) два, поэтому вспомогательная регрессия $\mathbf{x}_2$ на константу
и $\mathbf{x}_3$ — парная, и её $R^2$ равен квадрату коэффициента корреляции
(п.~\ref{sec:mc}). Средние: $\overline{x}_2 = 9$, $\overline{x}_3 = 2$.
Отклонения:
\begin{center}
\begin{tabular}{crrrr}
\toprule
 & 1 & 2 & 3 & 4 \\
\midrule
$x_2 - \overline{x}_2$ & $1$ & $-3$ & $3$ & $-1$ \\
$x_3 - \overline{x}_3$ & $1$ & $-1$ & $0$ & $0$ \\
\bottomrule
\end{tabular}
\end{center}
\[
  \sum (x_2 - \overline{x}_2)^2 = 20,
  \qquad
  \sum (x_3 - \overline{x}_3)^2 = 2,
  \qquad
  \sum (x_2 - \overline{x}_2)(x_3 - \overline{x}_3) = 1 + 3 = 4,
\]
\[
  R_2^2 = R_3^2 = r^2(x_2, x_3) = \frac{4^2}{20\cdot 2} = 0{,}4 .
\]
По \eqref{eq:vif}
\[
  \mathrm{VIF}_2 = \mathrm{VIF}_3 = \frac{1}{1 - 0{,}4} = \frac53 \approx 1{,}67 .
\]
Это далеко от порога 10: при $a = 2$ мультиколлинеарность слабая. Для
свободного члена VIF обычно не считают: вспомогательная регрессия
«константы на остальные регрессоры с константой» не имеет смысла.
\end{solution}
